 \Def Область целостности $K$ называется \textit{нетёровым кольцом}, если выполнено одно из следующих условий:
\begin{itemize}
    \item [1] Если $I_0 \subset I_1 \subset I_2 \subset . . . $ -- последовательность идеалов в $K$, то $\exists \ m \in \N : \forall \ n \geqslant m :I_n = I_m$.
    \item [2] Любой идеал в $K$ является конечнопорожденным.
    \item[3] Если $(a_0) \subset (a_0, a_1) \subset (a_0, a_1, a_2) ... $ - последовательность идеалов в $K$, то $\exists m \in \mathbb{N} \ \forall \ n \geq m: \ a_n \in (a_1, ..., a_m)$
\end{itemize}
\Proof
 \begin{itemize}
     \item []$(1 \Longrightarrow 3)$  Тривиально, это частный случай более общего утверждения.
       \item []$(3 \Longrightarrow 2)$ Пусть $I$ - ненулевой идеал в $K$. Выберем $a_0 \in I$. \newline Если $I = (a_0)$, то получено требуемое. \newline Иначе -- выберем $a_1 \in I\backslash (a_0)$. Если $I = (a_0, a_1)$, то получено требуемое. Иначе -- выберем $a_2 \in I\backslash (a_0, a_1)$ и продолжим процесс. Предположим, что этот процесс не завершается, тогда в $I$ есть нестабилизирующаяся последовательность идеалов $(a_0) \subset (a_0, a_1) \subset (a_0, a_1, a_2) \subset ...$ , что невозможно.
       \item []$(2 \Longrightarrow 1)$ Пусть $I_0 \subset I_1 \subset I_2 \subset . . .$ -- последовательность идеалов в $K$. Как уже было
доказано, $I :=  \bigcup\limits_{j=0}^{\infty} I_j$ - тоже идеал (билет 4.22), тогда $I$ конечнопорожден, то есть $I = (a_1, . . . , a_n)$
для некоторых $a_1, . . . , a_n \in K$. Но тогда $\exists m := $min$\{j \in \Z_{\geqslant 0} : a_1, . . . , a_n \in I_j\} :
\forall n \geq m : I_n = I_m.$
 \end{itemize}
 \EndProof